\documentclass[11pt]{article}
\usepackage[margin=0.75in]{geometry}
\usepackage{booktabs,hyperref}
\hypersetup{colorlinks=true,linkcolor=blue,urlcolor=blue}
\setlength{\parindent}{0pt}
\setlength{\parskip}{5pt}
\emergencystretch=2em
\newcommand{\heading}[1]{\section*{#1}}
\title{Impala acceleration: results and lessons}
\author{Ayan Biswas\quad ayanju04@gmail.com}
\date{September 30, 2026\quad---\quad Steps 1--9}
\begin{document}
\maketitle
\heading{What we achieved}
For the tested pooled copper calibration, the final experimental GPU workflow
completed four chains in \textbf{154.17 seconds}, compared with
\textbf{205.90 seconds on the optimized parallel CPU} and a historical
\textbf{962.03-second original sequential reference}.

\begin{center}
\begin{tabular}{lrr}
\toprule
Implementation & Sampling time & Relative throughput\\
\midrule
Original CPU, four sequential chains & 962.03 s & 1.00$\times$\\
Optimized CPU, four parallel chains & 205.90 s & 4.67$\times$\\
Optimized GPU, four coordinated chains & 154.17 s & 6.24$\times$\\
\bottomrule
\end{tabular}
\end{center}
\textbf{The fair incremental GPU comparison is 1.34$\times$ the optimized CPU
throughput}, or about 25\% less elapsed time. The 6.24$\times$ historical gain
combines CPU optimizations, parallel chains, GPU execution and worker improvements;
it must not be described as a GPU-only speedup.

\heading{Benchmark and comparison limits}
The long benchmark uses four independent seeds (42, 123, 456, 789), 40,000 stored
states per chain and 50 temperatures, with the existing pooled Cu/PTW example.
Each chain has its own random-number stream and adaptation state. Only forward-model
requests are combined for GPU evaluation.

The hardware is a Core Ultra 9 285HX CPU and RTX PRO 5000 Blackwell laptop GPU.
One numerical-library thread is used per worker. Timings exclude startup/JIT,
saving and diagnostics, while including initialization inside sampling, cache
misses, CPU work, transfers and coordination. The original reference is a historical
sum of sequential chain times; current CPU/GPU timings are group wall times.

The final CPU/GPU figures come from the same step-8 comparison. Each long backend
was measured once; repeated short controls provide supporting evidence. Results
are specific to this workload and machine. They are not universal performance
claims or uncertainty bounds.

\heading{Practical status}
The CPU implementation remains available. GPU support is an \textbf{optional,
experimental benchmark path for the tested Cu model}, not a finished general-purpose
production backend. GPU adapter cleanup, automatic backend selection/failover and
Mac validation remain future work. Sampler distributions, proposal order and
adaptation rules were preserved. Hierarchical mixing was not changed.

\newpage
\heading{How the improvements were achieved}
\renewcommand{\arraystretch}{1.2}
\begin{tabular}{@{}p{0.25\linewidth}p{0.70\linewidth}@{}}
\toprule
Change & Why it helped and what the measurements showed\\
\midrule
Cache fixed preparation & Reused experimental geometry and history inputs instead
of rebuilding them. Modest early gains; original fixture outputs stayed exact.\\
Compile CPU history loop & Used an optional Numba driver for the repeated material
history loop. Pooled short-run sampling was about 1.97--1.98$\times$ faster than
the cached Python path in the step-2 tests.\\
Batch independent chains & Combined separate chains' forward-model requests into
larger GPU calls, preserving each chain's RNG and sampler state. Batching alone
was insufficient: the first version still lost to parallel CPU.\\
Keep geometry on the GPU & Cached device geometry, performed interpolation on the
GPU and returned only needed predictions. This was the first coordinated GPU
version to beat the parallel CPU in the long benchmark.\\
Tune history kernel & Selected 32-thread blocks and moved invariant calculations
outside the integration loop, preserving float64 arithmetic.\\
Batch Gibbs bookkeeping & Vectorized repeated temperature-wise noise and likelihood
calculations on CPU while retaining draw order and RNG state. This reduced the
CPU work feeding the GPU.\\
Batch proposal likelihoods & Reduced scalar likelihood calls, but delivered only
a small additional GPU gain. CPU timing evidence was mixed; kept opt-in.\\
\bottomrule
\end{tabular}

\heading{Measured milestones: fresh controls within each step}
\begin{center}
\begin{tabular}{llrr}
\toprule
Step & Change & Previous GPU & New GPU\\
\midrule
5 & Resident geometry/interpolation & -- & 214.32 s\\
6 & Kernel tuning & 213.56 s & 202.71 s\\
7 & Batched Gibbs bookkeeping & 196.58 s & 157.59 s\\
8 & Batched proposal likelihoods & 156.06 s & 154.17 s\\
\bottomrule
\end{tabular}
\end{center}
Step 5 beat its fresh parallel CPU control (252.16 s) by 17.7\% in throughput.
Step 6 improved GPU throughput by 5.4\%; step 7 by 24.7\%; step 8 by 1.2\%.
The step-8 short median gain was 3.0\%, and CPU results were inconsistent.
\textbf{Do not multiply these percentages}: each step remeasured its controls,
and operating conditions/timing variation changed between runs.

\heading{What did not work well}
The first history-only GPU implementation was slower end-to-end than compiled CPU
for a single chain (211.88 versus 186.68 s). Initial four-chain GPU batching was
also slower than parallel CPU (334.53 versus 252.64 s). Reducing host work and
transfers mattered more than merely moving the history calculation to a GPU.
The final small proposal-batching gain does not establish a robust general speedup.

\newpage
\heading{How stochastic accuracy was checked}
We measured \textbf{effective samples per second (ESS/s)}, not just iterations
per second. Faster iterations are useful only if posterior quality and sample
correlation remain acceptable.

Validation combined an unchanged 256-case original prediction/likelihood fixture,
seeded comparisons against frozen original CPU traces, full-temperature short-run
checks and four-chain long-run posterior diagnostics. The long diagnostics discard
20,000 states per chain, without thinning, and cover 16 parameter/noise quantities.

\begin{center}
\begin{tabular}{lrr}
\toprule
Diagnostic & Required & Final GPU result\\
\midrule
Maximum R-hat & $<1.01$ & 1.002585\\
Minimum bulk ESS & $\geq400$ & 2253.48\\
Minimum tail ESS & $\geq400$ & 3345.25\\
Minimum bulk ESS/s & Comparative metric & 14.62\\
\bottomrule
\end{tabular}
\end{center}
The optimized CPU delivered 10.94 minimum bulk ESS/s in the same step-8 study,
so the GPU improvement was about 34\%. MCSE/SD, posterior-mean differences and
prediction-mean differences also passed the predeclared gates.

Optimized CPU parameter/noise/counter/prediction arrays reproduced the frozen
original CPU traces exactly. GPU parameter draws and counters matched CPU exactly;
noise and prediction differences were at floating-point rounding scale and passed
unchanged tolerances. Later GPU changes preserved saved GPU outputs bit-for-bit,
so their measured throughput gains also improved ESS/s for those runs.

The latest complete hardware-enabled regression suite passed \textbf{135 tests}.
Step 9 then passed 19 focused instrumentation/GPU tests, and all saved arrays
matched across 12 successful profiling runs. These are separate test runs, not
154 distinct tests. Bitwise agreement is evidence for these tested environments
and seeds, not a promise across hardware or numerical-library versions.

\heading{Why we stopped GPU optimization here}
Step 9 separated overlapping timers. In a representative standard instrumented
run, GPU kernels accounted for 49.8\% of elapsed time, other model work 17.8\%,
waiting for worker requests 24.3\%, and request decoding/reply sending 6.6\%.
Assembly and remaining bookkeeping accounted for about 1.6\%.
Polling includes necessary sampler work and cannot simply be removed. Detailed
instrumentation itself increased median runtime by about 2.2\%.

Larger batches revealed geometry preparation costs, mostly first-use setup rather
than cache evictions. These findings identify future targets, but do not establish
an asynchronous scheduling speedup. Given diminishing incremental returns, it is
reasonable to move to other acceleration work. \textbf{Hierarchical mixing remains
unresolved} and should be studied separately because it may require algorithm changes.

\heading{Where to find the evidence}
\small\raggedright
Checkpoint: \texttt{perf-step9-20260930} (\texttt{33e7da6}); branch
\texttt{perf/gpu-coordination-profile}. See \href{CHANGELOG.md}{the performance change log},
\texttt{STEP1.md}--\texttt{STEP9.md} and \path{results/step*-results.json} for commands,
measurements, diagnostic CSVs and SHA-256 manifests. Raw artifacts reside in workspace
\path{benchmark/changes/} and \path{benchmark/results/reference/}.
This consolidated PDF and its LaTeX source are included in the repository.
\end{document}
